
% ============================================================
\section{Introduzione e motivazione}
% ============================================================

L'imaging non invasivo di cellule biologiche viventi nel loro ambiente naturale è un
problema centrale nella biofisica moderna. Le tecniche di microscopia classica a
contrasto di fase (Zernike) o a contrasto di interferenza differenziale (Nomarski) sono
strumenti preziosi per rendere visibili campioni trasparenti e non colorati, ma offrono
soltanto informazioni qualitative: non è possibile, con esse, ricavare in modo assoluto
il ritardo di fase prodotto dal campione.
% ? sistemare questa parte su in base al campione dato
% * riscrivere le parole in italiano tecnice, in inglese

Le tecniche interferometriche digitali superano questo limite fornendo una misura
\emph{quantitativa} dello sfasamento del fronte d'onda trasmesso attraverso il
campione. In particolare, la \textbf{Microscopia Olografica Digitale} (DHM,
\textit{Digital Holographic Microscopy}) consente di ricostruire numericamente, a
partire da un singolo ologramma registrato da una camera CCD, sia l'immagine di
ampiezza che quella di \emph{contrasto di fase quantitativo} con risoluzione
interferometrica, senza l'impiego di agenti di contrasto.

Il vantaggio fondamentale è che la DHM fornisce una mappa spaziale del
\textbf{cammino ottico} (OPL, \textit{Optical Path Length}) punto per punto nel piano
del campione. Come si vedrà, tale grandezza dipende in modo \emph{accoppiato} da
due proprietà fisiche della cellula:
\begin{enumerate}
  \item lo \textbf{spessore locale} $h$ della cellula (morfometria), e
  \item l'\textbf{indice di rifrazione integrale} $\bar{n}_c$, ovvero la media
        dell'indice di rifrazione intracellulare lungo lo spessore.
\end{enumerate}

Misurare queste due quantità separatamente da un'unica immagine di fase è un
problema sottodeterminato. Questa relazione presenta, tra le altre cose, una procedura sperimentale
originale — la \textbf{procedura di disaccoppiamento} — che risolve tale
indeterminazione in modo elegante e rigoroso. % ? basato su articolo ...

% ============================================================
\section{Principio fisico: cammino ottico e sfasamento}
% ============================================================

\subsection{L'indice di rifrazione e la velocità della luce}

L'indice di rifrazione di un mezzo è definito come
\[
  n = \frac{c}{v},
\]
dove $c$ è la velocità della luce nel vuoto e $v$ è la velocità di fase nel mezzo.
%Materiali biologici come il citoplasma cellulare hanno indici di rifrazione
%superiori a quello del mezzo extracellulare acquoso ($n_m \approx 1.33$--$1.34$),
%risultando in un rallentamento della luce che attraversa la cellula rispetto a quella
%che viaggia nel mezzo circostante.
Materiali o soluzioni che simulano il citoplasma cellulare hanno indici di rifrazione
superiori a quello del mezzo extracellulare acquoso ($n_m \approx 1.33$--$1.34$) (dato in questa anologia dal campione sintetico di nastro adesivo)
risultando in un rallentamento della luce che attraversa il campione rispetto a quella
che viaggia nel mezzo circostante.

\subsection{Definizione del cammino ottico}

Sia $z$ la coordinata assiale (lungo l'asse ottico del microscopio). Il
\textbf{cammino ottico} $\mathrm{OPL}$ di un raggio che percorre un mezzo
\textit{stratificato} di spessore totale $D$ è
\[
  \mathrm{OPL} = \int_0^D n(z)\,dz,
\]
dove $n(z)$ è il profilo dell'indice di rifrazione lungo $z$.

\subsection{Configurazione sperimentale della DHM}

Il campione (un multistrato di layer di nastro adesivo) è posto in una camera di perfusione (alloggiamento circolare apposito) di
altezza $D$ riempita di soluzione con indice di rifrazione $n_m$ uniforme e costante (SOLUZIONE DI SACCAROSO A PERCENTUALE, CON UN n MISURATO).
Il campione occupa localmente una regione di spessore $h_i$ (per ciascun pixel $i$
dell'immagine), con \textbf{indice di rifrazione interno} al campione $n_{c,i}(z)$ variabile in senso
assiale.

Il cammino ottico per il pixel $i$ è:
\begin{equation}
  \mathrm{OPL}_i
  = \int_0^{h_i} n_{c,i}(z)\,dz + n_m(D - h_i)
  = \left(\bar{n}_{c,i} - n_m\right)h_i + n_m D,
  \label{eq:OPL}
\end{equation}
dove si è definito l'\textbf{indice di rifrazione integrale} come la media integrale:
\[
  \bar{n}_{c,i} = \frac{1}{h_i}\int_0^{h_i} n_{c,i}(z)\,dz.
\]

\subsection{Relazione tra OPL e fase misurata}

La fase $\varphi_i$ misurata dal DHM per il pixel $i$ è legata all'OPL dalla
relazione:
\[
  \mathrm{OPL}_i = \frac{\lambda}{2\pi}\,\varphi_i,
\]
dove $\lambda$ è la lunghezza d'onda del laser. Sostituendo~\eqref{eq:OPL}:
\begin{equation}
    \color{red}
  \varphi_i = \frac{2\pi}{\lambda}\left[\left(\bar{n}_{c,i} - n_m\right)h_i + n_m D\right].
  \label{eq:fase_completa}
\end{equation}

Il termine $n_m D$ è una costante di riferimento misurabile in zone prive di campione (o di cellule)
e viene sottratta sperimentalmente. Il segnale \emph{di fase} è quindi:
\begin{equation}
  \color{red}
  \varphi_i = \frac{2\pi}{\lambda}\left(\bar{n}_{c,i} - n_m\right)h_i.
  \label{eq:fase_cella}
\end{equation}

\paragraph{Il problema dell'accoppiamento.}
L'equazione~\eqref{eq:fase_cella} contiene \textbf{due incognite} per ciascun
pixel $i$: lo spessore $h_i$ e l'indice di rifrazione integrale
$\bar{n}_{c,i}$. Una singola misura di fase non è sufficiente a determinare
entrambe le quantità. Si noti in particolare che variazioni di $h_i$ e di
$\bar{n}_{c,i}$ possono produrre identiche variazioni di $\varphi_i$, rendendo
ambigua l'interpretazione del segnale di fase da sola.

% ============================================================
\section{Stabilità del sistema e prestazioni della DHM}
% ============================================================

Prima di illustrare la procedura di disaccoppiamento è importante quantificare le
prestazioni del sistema, poiché esse determinano la precisione finale delle misure.

\begin{itemize}
  \item \textbf{Stabilità temporale della fase:} mediando 16 immagini di fase
        consecutive, si ottiene una stabilità temporale equivalente a
        $\lambda/1800 \approx 0{,}2°$ per ciascun pixel.
  \item \textbf{Sensibilità verticale (assiale):} la stabilità di fase di $0{,}2°$
        corrisponde, per campioni con indice di rifrazione tipico
        $\bar{n}_c \approx 1{,}375$, a una sensibilità verticale di circa
        \textbf{11 nm}.
  \item \textbf{Risoluzione laterale:} limitata dalla diffrazione,
        $\approx 0{,}6\ \mu\mathrm{m}$ per un obiettivo con
        $\mathrm{NA} = 0{,}8$ a $\lambda = 658\ \mathrm{nm}$ (criterio di
        Rayleigh per illuminazione coerente).
  \item \textbf{Tempo di acquisizione:} limitato dal tempo di esposizione della
        camera ($\sim 20\ \mu\mathrm{s}$), con ricostruzione in tempo reale a
        più di 10 immagini/s.
  \item \textbf{Bassa intensità di illuminazione:} l'irradianza sul campione è
        $\sim 200\ \mu\mathrm{W/cm^2}$, svariati ordini di grandezza inferiore a
        quella tipica della microscopia confocale a scansione laser.
\end{itemize}

% ============================================================
\section{La procedura di disaccoppiamento}
\label{sec:decoupling}
% ============================================================

\subsection{Idea fondamentale}

Per risolvere l'indeterminazione dell'equazione~\eqref{eq:fase_cella}, si introduce
una \emph{seconda misura di fase indipendente} sullo \emph{stesso campione},
ottenuta modificando in modo controllato e noto l'indice di rifrazione del mezzo
extracellulare (mezzo fuori dal campione). In questo modo si ottiene un \textbf{sistema di due equazioni in due
incognite} ($h_i$ e $\bar{n}_{c,i}$) che può essere risolto analiticamente.

\subsection{Protocollo sperimentale}

La procedura prevede la perfusione sequenziale di \textbf{due soluzioni}:

\begin{enumerate}
  \item \textbf{Soluzione standard}: mezzo extracellulare (acqua + saccarosio) con indice di rifrazione
        $n_m$ e osmolarità $\Pi_0$. Si registra l'ologramma e si ricostruisce la
        mappa di fase $\varphi_{1,i}$.

  \item \textbf{Soluzione di disaccoppiamento}: mezzo extracellulare con indice di
        rifrazione $n'_m = n_m + \delta n$. Si registra il secondo ologramma e si ottiene $\varphi_{2,i}$.
\end{enumerate}

Il requisito di \textit{isosmoticità} è cruciale: esso assicura che $h_i$ non cambi tra le
due misure, in modo che l'unica differenza tra $\varphi_{1,i}$ e $\varphi_{2,i}$
sia dovuta esclusivamente alla variazione $\delta n$ dell'indice del mezzo
extracellulare.

\paragraph{Come si ottiene $\delta n$:}
Si sostituisce una molecola idrofila osmotica (presente nella soluzione standard
come agente osmotico) con una quantità equimolare di una diversa molecola idrofila
avente \emph{maggiore polarizzabilità molare}. Poiché la concentrazione molare
totale rimane invariata, l'osmolarità è conservata, mentre l'indice di rifrazione
aumenta di $\delta n$ noto e misurabile.

\subsection{Il sistema di equazioni}

Con le convenzioni sopra introdotte, le due misure di fase sono:

\begin{align}
  \varphi_{1,i} &= \frac{2\pi}{\lambda}\left(\bar{n}_{c,i} - n_m\right)h_i,
  \label{eq:phi1} \\[6pt]
  \varphi_{2,i} &= \frac{2\pi}{\lambda}\left(\bar{n}_{c,i} - n_m - \delta n\right)h_i.
  \label{eq:phi2}
\end{align}

Sottraendo~\eqref{eq:phi2} da~\eqref{eq:phi1}:
\begin{equation}
  \varphi_{1,i} - \varphi_{2,i}
  = \frac{2\pi}{\lambda}\,\delta n\, h_i.
  \label{eq:differenza}
\end{equation}

\subsection{Derivazione della formula per lo spessore cellulare}
\label{subsec:spessore}

Dall'equazione~\eqref{eq:differenza} si ricava immediatamente lo spessore:

\begin{equation}
  \boxed{
    h_i = \frac{\lambda}{2\pi}\,\frac{\varphi_{1,i} - \varphi_{2,i}}{\delta n}
  }
  \label{eq:spessore}
\end{equation}

\paragraph{Interpretazione.} Lo spessore $h_i$ è determinato dalla
\emph{differenza} tra i due segnali di fase, normalizzata per la variazione
$\delta n$ introdotta nella soluzione. È notevole il fatto che il risultato non
dipenda dall'indice di rifrazione intracellulare $\bar{n}_{c,i}$, che è
sconosciuto: il contributo di quest'ultimo si cancella nella differenza
$\varphi_{1,i} - \varphi_{2,i}$. Ciò rende la misura di $h_i$ \emph{intrinsecamente
robusta} rispetto all'eterogeneità dell'intracellulare.

Notiamo anche la consistenza dimensionale: $[\lambda/(2\pi)]$ ha le dimensioni di
una lunghezza, $[\delta n]$ è adimensionale (variazione dell'indice), e
$[\varphi_{1,i} - \varphi_{2,i}]$ è in radianti (adimensionale), quindi $[h_i]$
risulta correttamente una lunghezza.

\subsection{Derivazione della formula per l'indice di rifrazione integrale}
\label{subsec:indice}

Per ricavare $\bar{n}_{c,i}$, si sostituisce~\eqref{eq:spessore}
in~\eqref{eq:phi1}:
\[
  \varphi_{1,i}
  = \frac{2\pi}{\lambda}\left(\bar{n}_{c,i} - n_m\right)
    \cdot \frac{\lambda}{2\pi}\,\frac{\varphi_{1,i} - \varphi_{2,i}}{\delta n}.
\]
Semplificando:
\[
  \varphi_{1,i}
  = \left(\bar{n}_{c,i} - n_m\right)
    \frac{\varphi_{1,i} - \varphi_{2,i}}{\delta n}.
\]
Risolvendo per $\bar{n}_{c,i}$:
\[
  \bar{n}_{c,i} - n_m
  = \frac{\delta n\,\varphi_{1,i}}{\varphi_{1,i} - \varphi_{2,i}},
\]
e quindi:
\begin{equation}
  \boxed{
    \bar{n}_{c,i}
    = \frac{\delta n\,\varphi_{1,i}}{\varphi_{1,i} - \varphi_{2,i}} + n_m
  }
  \label{eq:indice}
\end{equation}

\paragraph{Interpretazione.}
La formula~\eqref{eq:indice} mostra che l'indice di rifrazione integrale è
determinato dal \emph{rapporto} tra il segnale di fase nella prima misura
$\varphi_{1,i}$ e la differenza $\varphi_{1,i} - \varphi_{2,i}$, scalato per
$\delta n$ e traslato di $n_m$. Si noti che:
\begin{itemize}
  \item Se $\delta n \to 0$ il metodo perde di sensibilità (il denominatore va
        a zero); occorre quindi una perturbazione $\delta n$ sufficientemente
        grande da essere misurabile, ma non così grande da perturbare
        biologicamente la cellula.
  \item $\bar{n}_{c,i}$ non dipende dalla lunghezza d'onda $\lambda$ né da $D$,
        essendo questi parametri già stati eliminati nel passaggio dalla fase al
        segnale cellulare~\eqref{eq:fase_cella}.
  \item Il risultato è valido pixel per pixel: si ottiene quindi una
        \emph{mappa spaziale} dell'indice di rifrazione integrale, non solo un
        valore medio.
\end{itemize}

\subsection{Riepilogo del sistema: forma matriciale}
\label{subsec:matriciale}

Il sistema di equazioni~\eqref{eq:phi1}--\eqref{eq:phi2} può essere scritto in
forma matriciale per evidenzirne la struttura lineare. Poniamo
$x_i = \bar{n}_{c,i}\,h_i$ e $y_i = h_i$. Allora:
\[
\begin{pmatrix} \varphi_{1,i} \\ \varphi_{2,i} \end{pmatrix}
= \frac{2\pi}{\lambda}
\begin{pmatrix} 1 & -n_m \\ 1 & -(n_m+\delta n) \end{pmatrix}
\begin{pmatrix} x_i \\ y_i \end{pmatrix}.
\]
Il determinante della matrice dei coefficienti è:
\[
  \Delta = \frac{(2\pi)^2}{\lambda^2}
  \left[-(n_m+\delta n) + n_m\right]
  = -\frac{(2\pi)^2}{\lambda^2}\,\delta n \neq 0
  \quad \text{se } \delta n \neq 0.
\]
La non-singolarità del sistema è garantita dal fatto che $\delta n \neq 0$,
confermando che la procedura di disaccoppiamento è ben posta se e solo se le due
soluzioni hanno indice di rifrazione \emph{diverso}.

% ============================================================
\section{Precisione e limiti della procedura}
% ============================================================

\subsection{Stime teoriche di accuratezza per campioni fissi}

Data la stabilità di fase di $\sigma_\varphi \approx 0{,}2°$ per il sistema DHM
considerato, le accuratezze teoriche delle quantità ricostruite per un singolo pixel
sono:
\begin{align*}
  \sigma_{h_i} &\approx \frac{\lambda}{2\pi}\,\frac{\sqrt{2}\,\sigma_\varphi}{\delta n}
               \approx 100\ \mathrm{nm}, \\[4pt]
  \sigma_{\bar{n}_{c,i}} &\approx 0{,}0004.
\end{align*}

Queste accuratezze sono ottenibili solo su \textbf{campioni fissi}, poiché
richiedono che il campione rimanga perfettamente immobile tra le due acquisizioni.

\subsection{Effetti dei micro-movimenti su campioni viventi}

Per campioni \emph{viventi}, la cellula esegue spontaneamente micro-movimenti
durante il tempo di scambio della soluzione nel canale di perfusione (tipicamente
$\sim 30\ \mathrm{s}$). Tali movimenti inducono fluttuazioni temporali del segnale
di fase dell'ordine di $0{,}5$--$1{,}5°$ per pixel, peggiorando significativamente
le accuratezze puntuali.

\subsection{La media spaziale come rimedio}

Per compensare gli artefatti da micro-movimenti, si applica una \textbf{media
spaziale} dell'indice di rifrazione su tutta la superficie cellulare:
\[
  \bar{n}_c = \frac{1}{N_\mathrm{cell}}\sum_{i \in S_\mathrm{cell}} \bar{n}_{c,i}.
\]
La superficie cellulare $S_\mathrm{cell}$ viene determinata tramite un algoritmo di
rilevazione dei bordi basato su gradienti, seguito da una procedura di erosione che
rimuove i pixel periferici con basso rapporto segnale-rumore.

Questo indice mediato $\bar{n}_c$ presenta una stabilità temporale di $0{,}0003$
su intervalli di 30 secondi. Come prima approssimazione, si utilizza poi $\bar{n}_c$
(invece di $\bar{n}_{c,i}$ pixel per pixel) nella formula dello spessore~\eqref{eq:spessore},
riscritta come:
\[
  h_i \approx \frac{\lambda}{2\pi}\,\frac{\varphi_{1,i}}{\bar{n}_c - n_m}.
\]

\paragraph{Accuratezze pratiche su campioni viventi.}
Con la media spaziale:
\begin{itemize}
  \item Lo spessore locale $h_i$ ha un'accuratezza di $\sim 1\ \mu\mathrm{m}$.
  \item Mediando su $\sim 36$ pixel adiacenti, l'accuratezza locale migliora a
        $\sim 0{,}7\ \mu\mathrm{m}$.
  \item Lo spessore medio di una regione cellulare ha un'accuratezza di poche
        decine di nanometri.
  \item La variazione spaziale dell'indice integrale è stimata in $\sim 0{,}005$.
\end{itemize}

% ============================================================
\section{Discussione critica delle formule (4) e (5)}
% ============================================================

\subsection{Formula (4): misura dell'indice di rifrazione integrale}

\begin{equation*}
  \bar{n}_{c,i}
  = \frac{\delta n\,\varphi_{1,i}}{\varphi_{1,i} - \varphi_{2,i}} + n_m
  \tag{4}
\end{equation*}

\paragraph{Struttura matematica.}
Il termine $\dfrac{\varphi_{1,i}}{\varphi_{1,i} - \varphi_{2,i}}$ è un rapporto
di fasi. Poiché $\varphi_{2,i} < \varphi_{1,i}$ (la seconda soluzione ha indice
più alto, quindi lo sfasamento prodotto dalla cellula rispetto al mezzo è minore),
il denominatore è positivo e l'indice risultante è maggiore di $n_m$, come
fisicamente richiesto.

\paragraph{Dipendenza da $\delta n$.}
La sensibilità della misura aumenta con $\delta n$: una perturbazione maggiore
produce una differenza $\varphi_{1,i} - \varphi_{2,i}$ più grande, riducendo
l'errore relativo sul rapporto. Tuttavia, $\delta n$ non può essere arbitrariamente
grande perché la molecola utilizzata per modificare l'indice non deve
attraversare la membrana cellulare (la si sceglie idrofila e di alto peso
molecolare) e l'isosmoticità deve essere mantenuta.

\paragraph{Limite per $\delta n \to \varphi_{1,i}$ (caso limite).}
Se $\varphi_{2,i} \to 0$, ovvero la seconda soluzione ha indice di rifrazione
pari a $\bar{n}_{c,i}$ stesso (la cellula diventa invisibile in fase), si ottiene
$\bar{n}_{c,i} = \delta n + n_m = n'_m$, il che è coerente.

\subsection{Formula (5): misura dello spessore cellulare}

\begin{equation*}
  h_i
  = \frac{\lambda}{2\pi}\,\frac{\varphi_{1,i} - \varphi_{2,i}}{\delta n}
  = \frac{\lambda}{2\pi}\,\frac{\varphi_{1,i}}{\bar{n}_{c,i} - n_m}
  \tag{5}
\end{equation*}

\paragraph{Due forme equivalenti.}
La formula (5) è presentata in due forme equivalenti:
\begin{enumerate}
  \item La \textbf{forma operativa} (a sinistra): utilizza direttamente la
        differenza delle fasi misurate e $\delta n$ noto. Non richiede la
        conoscenza di $\bar{n}_{c,i}$.
  \item La \textbf{forma esplicita} (a destra): si ottiene sostituendo la
        definizione di $\bar{n}_{c,i}$ e mostra che $h_i$ è legato alla fase
        nella prima soluzione, normalizzata per il contrasto di indice
        $\bar{n}_{c,i} - n_m$.
\end{enumerate}

\paragraph{Robustezza rispetto all'indice intracellulare.}
La forma operativa della (5) mostra che lo spessore può essere misurato
\emph{senza conoscere a priori} l'indice di rifrazione intracellulare. Questo è
il punto di forza del metodo: la misura morfometrica è separata e indipendente
dalla misura dell'indice.

\paragraph{Consistenza con la fisica.}
Nella forma esplicita, l'espressione $\dfrac{\lambda}{2\pi}\dfrac{\varphi_{1,i}}
{\bar{n}_{c,i}-n_m}$ ha un'interpretazione fisica immediata: il numeratore
$\dfrac{\lambda}{2\pi}\varphi_{1,i}$ è il cammino ottico extra prodotto dalla
cellula rispetto al mezzo, mentre il denominatore $\bar{n}_{c,i} - n_m$ è
l'eccesso di indice di rifrazione; il loro rapporto fornisce direttamente
lo spessore fisico.

\subsection{Relazione tra le due formule}

Si noti che le formule (4) e (5) non sono indipendenti: derivano dalla risoluzione
dello stesso sistema lineare~\eqref{eq:phi1}--\eqref{eq:phi2}. La coerenza tra
le due può essere verificata sostituendo~\eqref{eq:spessore}
in~\eqref{eq:fase_cella} e recuperando~\eqref{eq:indice}. Il sistema è dunque
\emph{completo e non ridondante}: la misura di due fasi $(\varphi_{1,i},
\varphi_{2,i})$ con un $\delta n$ noto determina univocamente sia $h_i$ che
$\bar{n}_{c,i}$.

% ============================================================
\section{Stima del volume cellulare dalla morfometria}
% ============================================================

Una volta noti gli spessori $\{h_i\}$ per tutti i pixel all'interno della superficie
cellulare $S_\mathrm{cell}$, il volume cellulare può essere stimato come:
\[
  V_\mathrm{cell}
  \approx N_\mathrm{cell}\,\frac{S_\mathrm{pixel}}{M^2}\,\bar{h}_\mathrm{cell},
\]
dove $N_\mathrm{cell}$ è il numero di pixel nella superficie cellulare,
$S_\mathrm{pixel}$ è l'area di un singolo pixel nella camera CCD, $M$ è l'ingrandimento
del DHM, e $\bar{h}_\mathrm{cell} = \dfrac{1}{N_\mathrm{cell}}\displaystyle\sum_{i \in S_\mathrm{cell}} h_i$
è lo spessore medio della cellula.

% ============================================================
\section{Interpretazione fisica dell'indice di rifrazione integrale}
% ============================================================

L'indice di rifrazione integrale $\bar{n}_c$ è strettamente legato alla
\textbf{composizione biochimica} del citoplasma. Nell'approssimazione di
campo medio di Lorentz-Lorenz, l'indice di rifrazione di una miscela acquosa di
biomolecole aumenta linearmente con la concentrazione di materiale non acquoso
(proteine, ioni, lipidi, ecc.):
\[
  \bar{n}_c \approx n_{\mathrm{H_2O}} + \alpha_\mathrm{spec}\,C_\mathrm{protein},
\]
dove $\alpha_\mathrm{spec}$ è l'incremento specifico di rifrazione per unità di
concentrazione (tipicamente $\approx 0{,}185\ \mathrm{ml/g}$ per le proteine) e
$C_\mathrm{protein}$ è la concentrazione proteica intracellulare.

Pertanto, una \emph{diminuzione} di $\bar{n}_c$ segnala una \textbf{diluizione
intracellulare} (ad es.\ ingresso di acqua), mentre un \emph{aumento} segnala
un incremento di concentrazione di materiale disciolto.

% ============================================================
\section{Conclusioni}
% ============================================================

La procedura di disaccoppiamento descritta in questo articolo rappresenta una
soluzione elegante a un problema fondamentale nella microscopia di fase
quantitativa: \emph{come separare il contributo morfometrico (spessore) da quello
ottico (indice di rifrazione) nel segnale di fase}.

Il metodo si basa su due principi fisici semplici:
\begin{enumerate}
  \item La fase prodotta da una cellula dipende \emph{linearmente} sia da $h_i$
        che da $\bar{n}_{c,i}$ (equazione~\eqref{eq:fase_cella}).
  \item Modificando l'indice del mezzo extracellulare di una quantità nota
        $\delta n$ a volume cellulare costante, si ottiene un secondo vincolo
        lineare che rende il sistema determinato.
\end{enumerate}

Le formule chiave:
\begin{align*}
  \bar{n}_{c,i} &= \frac{\delta n\,\varphi_{1,i}}{\varphi_{1,i}-\varphi_{2,i}} + n_m,
  \tag{4} \\[6pt]
  h_i &= \frac{\lambda}{2\pi}\,\frac{\varphi_{1,i}-\varphi_{2,i}}{\delta n},
  \tag{5}
\end{align*}
sono entrambe derivabili in due passaggi algebrici elementari dalla differenza e
dal rapporto delle equazioni~\eqref{eq:phi1} e~\eqref{eq:phi2}.

Le prestazioni del DHM — stabilità di fase equivalente a $\lambda/1800$, risoluzione
laterale diffrattiva, tempo di acquisizione di $\sim 20\ \mu\mathrm{s}$, assenza di
agenti di contrasto — rendono questa tecnica particolarmente adatta allo studio
\emph{in tempo reale} della dinamica cellulare, con applicazioni che spaziano dalla
misurazione dello stress osmotico alla fisiopatologia cellulare, dalla biofisica delle
membrane alla valutazione farmacologica.

% ! ---------------------------------------------

% ================================================================
\section{Contesto e obiettivo dell'esperimento}
% ================================================================

Il presente documento descrive un esperimento di calibrazione e validazione della
procedura DHM, ispirato al metodo di disaccoppiamento di Rappaz \textit{et al.}
(Optics Express, 2005), ma adattato a un campione solido, fisso e con spessore
geometricamente noto.

\subsection{Il campione}

Il campione è costruito sovrapponendo $N = 8$ strati di nastro adesivo, il cui
spessore per singolo strato è stato misurato con un calibro:
\[
  h_\text{layer} \approx 30\ \mu\text{m}.
\]

\begin{quote}
\textcolor{red!70!black}{\textbf{Attenzione:}} \textbf{Incertezza sullo spessore.} La misura con calibro su materiale
flessibile ha un'incertezza elevata: stimare $\sigma_{h} \approx \pm 5\ \mu\text{m}$
è conservativo ma realistico. Questa incertezza si propaga su tutte le stime
quantitative derivate. Se disponibile, una misura indipendente con profilometro
o microscopia confocale aumenterebbe significativamente la precisione.
\end{quote}

Il campione complessivo (8 strati, $H = 8 \times h_\text{layer} \approx 240\ \mu\text{m}$)
è stato tagliato con un bisturi in modo da esporre i singoli strati come una
\textbf{scalinata}: ogni gradino $n$ ha spessore nominale
\[
  h_n = n \cdot h_\text{layer}, \qquad n = 1, 2, \dots, 8.
\]
Questa geometria fornisce una \textbf{verità di riferimento indipendente dalla DHM}:
se il metodo è corretto, le misure DHM devono riprodurre la sequenza
$h_1, h_2, \dots, h_8$ con incrementi costanti.

\subsection{Il mezzo esterno}

Il campione è immerso in soluzioni acquose di saccarosio a concentrazioni
nominali $c = 40\%, 60\%, 80\%$ (e una quarta concentrazione, vedere Nota~\ref{nota:100}).
Ogni soluzione ha un indice di rifrazione $n_m^{(k)}$ misurato con un
rifrattometro Abbe alla lunghezza d'onda del laser DHM e alla temperatura
ambiente al momento dell'acquisizione.

\begin{quote}
\label{nota:100}
\textcolor{red!70!black}{\textbf{Attenzione:}} \textbf{Nota sulla concentrazione ``100\%''.}
Saccarosio al 100\% in peso è un solido cristallino, non una soluzione.
Se la quarta concentrazione indicata è 100\%, si tratta probabilmente di:
(a) acqua pura (0\%, usata come riferimento), oppure
(b) una concentrazione in unità diverse (es.\ g/L).
Questa ambiguità deve essere risolta prima dell'analisi quantitativa.
Nel seguito il quarto punto è indicato come $c_4$ con $n_m^{(4)}$ da misurare.
\end{quote}

\subsection{Differenza fondamentale rispetto all'esperimento biologico}

Nel paper di riferimento il campione (una cellula vivente) ha spessore e indice
di rifrazione entrambi \textit{incogniti} e può muoversi tra le due acquisizioni.
Nel nostro caso:
\begin{itemize}
  \item Il campione è \textbf{rigido e fisso}: $h_n$ non cambia tra acquisizioni.
  \item Lo spessore $h_n$ è \textbf{noto a priori} dalla costruzione del campione.
  \item L'indice $n_\text{tape}$ è \textbf{uniforme} su tutto il campione (omogeneo). %(uniforme non proprio, ma comunque noto).
  \item L'unica incognita ottica è $n_\text{tape}$.
\end{itemize}
Questo riduce il problema da due incognite per pixel a \textbf{una sola incognita
globale}.

% ================================================================
\section{Principio fisico e equazione di misura}
% ================================================================

\subsection{Cammino ottico e fase DHM}

In un DHM a trasmissione, la fase $\varphi$ misurata per un campione di spessore
$h$ e indice di rifrazione $n_\text{tape}$, immerso in un mezzo con indice
$n_m$ (soluzione di saccarosio), è:

\begin{equation}
  \varphi^{(k)} = \frac{2\pi}{\lambda}
  \bigl(n_\text{tape} - n_m^{(k)}\bigr)\cdot h
  \label{eq:fase}
\end{equation}

dove:
\begin{itemize}
  \item $\lambda$ è la lunghezza d'onda del laser DHM
        (\textbf{dovrebbe essere a 635 nm}: usare il valore nominale del sistema);
  \item $n_m^{(k)}$ è l'indice della soluzione $k$, misurato col rifrattometro;
  \item $\varphi^{(k)}$ è la \textbf{fase unwrapped} mediata sul gradino
        (vedere Sezione~\ref{sec:unwrapping}).
\end{itemize}

Il termine di riferimento $n_m D$ (dove $D$ è l'altezza della camera) viene
eliminato sottraendo la fase misurata in una regione senza campione
(\textit{background subtraction}).

\subsection{La fase come funzione lineare di $n_m^{(k)}$}

Per un gradino fissato con spessore $h_n = n\,h_\text{layer}$, l'equazione
\eqref{eq:fase} è \textbf{lineare in $n_m^{(k)}$}:

\begin{equation}
  \varphi_n^{(k)}
  = \underbrace{-\frac{2\pi\, n\, h_\text{layer}}{\lambda}}_{=\,m_n} \cdot n_m^{(k)}
  + \underbrace{\frac{2\pi\, n\, h_\text{layer}}{\lambda} \cdot n_\text{tape}}_{=\,q_n}
  \label{eq:retta}
\end{equation}

Il grafico $\varphi_n^{(k)}$ in funzione di $n_m^{(k)}$ deve essere una
\textbf{retta} con:

\begin{center}
\begin{tabular}{lll}
\hline\hline
Parametro & Espressione & Informazione estratta \\
\hline
Pendenza $m_n$ & $-\dfrac{2\pi\, n\, h_\text{layer}}{\lambda}$ &
  Misura di $h_\text{layer}$ (se $\lambda$ noto) \\[10pt]
Zero della retta & $n_m^* = n_\text{tape}$ &
  Misura diretta di $n_\text{tape}$ \\[6pt]
Intercetta $q_n$ & $m_n \cdot n_\text{tape}$ &
  Ridondante (consistenza) \\
\hline\hline
\end{tabular}
\end{center}

\paragraph{Proprietà chiave: le rette di gradini diversi si annullano nello stesso punto.}
Il punto $n_m^* = n_\text{tape}$ non dipende da $n$: tutte le rette relative
a gradini diversi si azzerano alla stessa concentrazione. Questo è un
\textbf{test di consistenza interna} molto robusto: se le rette non convergono
in un unico punto, vi è un errore sistematico (drift termico, errore nella
maschera, unwrapping errato).

Fisicamente, $n_m^{(k)} = n_\text{tape}$ è il punto di \textbf{matching dell'indice}:
la soluzione esterna ha lo stesso indice del nastro, il campione diventa
otticamente invisibile e la fase si annulla.

Per soluzioni di saccarosio, l'indice di rifrazione a temperatura ambiente
e $\lambda \approx 650\ \text{nm}$ vale approssimativamente:

\begin{center}
\begin{tabular}{cc}
\hline\hline
Concentrazione w/w (\%) & $n_m$ (approssimativo) \\
\hline
0  (acqua pura)   & 1.333 \\
40                & 1.400 \\
60                & 1.442 \\
80                & 1.490 \\
\hline\hline
\end{tabular}
\end{center}

I valori esatti vanno sostituiti con quelli misurati dal rifrattometro del
laboratorio.

Se $n_\text{tape}$ è nell'intervallo tipico dei nastri adesivi in
polipropilene o PVC ($1.47$--$1.50$), il punto di matching si trova vicino
alla soluzione all'80\%, rendendo quella misura particolarmente informativa.

% ================================================================
\section{Problema critico: il phase wrapping}
\label{sec:unwrapping}
% ================================================================

\begin{quote}
\textcolor{red!70!black}{\textbf{Attenzione:}} \textbf{Questo è il punto più critico dell'intera analisi.}
Con i parametri del vostro esperimento la fase è tipicamente avvolta
(\textit{wrapped}) più volte. Usare la fase wrapped nelle formule
produce risultati privi di significato fisico.
\end{quote}

\subsection{Stima quantitativa del numero di avvolgimenti}

Per $h_\text{layer} = 30\ \mu\text{m}$, $\lambda = 650\ \text{nm}$ (valore
tipico DHM) e $\Delta n = n_\text{tape} - n_m \approx 0.08$
(nastro in soluzione al 40\% di saccarosio):

\[
  \varphi = \frac{2\pi}{\lambda}\,\Delta n \cdot h_\text{layer}
  = \frac{2\pi}{650 \times 10^{-9}} \times 0.08 \times 30 \times 10^{-6}
  \approx 23\ \text{rad} \approx 3.7 \times 2\pi.
\]

Questo significa quasi \textbf{4 giri completi}: la fase wrapped è un numero
tra $-\pi$ e $\pi$ che non porta nessuna informazione utile senza correzione.

Il numero di avvolgimenti al variare della concentrazione (e quindi di $\Delta n$)
per il gradino $n=1$:

\begin{center}
\begin{tabular}{ccccc}
\hline\hline
Concentrazione (\%) & $n_m$ & $\Delta n$ & $\varphi$ (rad) & N.\ giri \\
\hline
40 & 1.400 & $\sim$0.08 & $\sim$23 & $\sim$3.7 \\
60 & 1.442 & $\sim$0.04 & $\sim$11 & $\sim$1.8 \\
80 & 1.490 & $\sim$0.00--0.01 & $\sim$0--3 & $\sim$0--0.5 \\
\hline\hline
\end{tabular}
\end{center}

\noindent (valori approssimativi basati su $n_\text{tape} \approx 1.49$,
$\lambda = 650\ \text{nm}$, $h = 30\ \mu\text{m}$.)

\subsection{Come verificare che l'unwrapping sia corretto}

\begin{enumerate}
  \item \textbf{Ispezione visiva della mappa di fase unwrapped}: la superficie
        del gradino deve apparire piatta e uniforme, senza salti bruschi di
        $2\pi \approx 6.28\ \text{rad}$. Salti residui indicano fallimento
        dell'unwrapping.

  \item \textbf{Test di monotonia}: all'aumentare di $n$ (numero del gradino)
        la fase media $\langle\varphi_n\rangle$ deve aumentare \textit{monotonicamente}
        (a parità di soluzione). Se un gradino dà un valore inferiore al precedente,
        c'è un salto di $2\pi$ non corretto.

  \item \textbf{Test di segno}: per $n_\text{tape} > n_m^{(k)}$ la fase deve essere
        \textit{positiva}. Per $n_\text{tape} < n_m^{(k)}$ deve essere \textit{negativa}.
        \color{red!70!black}{Per la soluzione all'80\%, se siete vicini al matching, la fase dovrebbe
        essere vicina a zero con segno determinato dall'ordinamento degli indici.}
\end{enumerate}

% ================================================================
\section{Procedura di disaccoppiamento: recupero delle formule (4) e (5)}
% ================================================================

Lavorando a coppie di soluzioni $(k_1, k_2)$ con $\delta n = n_m^{(k_1)} - n_m^{(k_2)}$,
le due equazioni di misura per un gradino di spessore $h$ sono:

\begin{align}
  \varphi^{(k_1)} &= \frac{2\pi}{\lambda}
    \bigl(n_\text{tape} - n_m^{(k_1)}\bigr)\,h,
  \label{eq:phi1} \\[6pt]
  \varphi^{(k_2)} &= \frac{2\pi}{\lambda}
    \bigl(n_\text{tape} - n_m^{(k_2)}\bigr)\,h.
  \label{eq:phi2}
\end{align}

\subsection{Derivazione della formula per lo spessore}

Sottraendo \eqref{eq:phi2} da \eqref{eq:phi1}:
\[
  \varphi^{(k_1)} - \varphi^{(k_2)}
  = \frac{2\pi}{\lambda}\,\delta n\, h,
\]
da cui:

\begin{equation}
  h = \frac{\lambda}{2\pi}\cdot
  \frac{\varphi^{(k_1)} - \varphi^{(k_2)}}{\delta n}
  \tag{5$'$}
  \label{eq:spessore}
\end{equation}

Lo spessore ricostruito $h$ \textbf{non dipende da $n_\text{tape}$}:
il contributo dell'indice del nastro si cancella nella differenza.

\subsection{Derivazione della formula per l'indice di rifrazione}

Dividendo \eqref{eq:phi1} per \eqref{eq:spessore} ed esplicitando $n_\text{tape}$:

\begin{equation}
  n_\text{tape} = \frac{\delta n\cdot\varphi^{(k_1)}}
  {\varphi^{(k_1)} - \varphi^{(k_2)}} + n_m^{(k_1)}
  \tag{4$'$}
  \label{eq:indice}
\end{equation}

Queste sono esattamente le formule (4) e (5) del paper di Rappaz \textit{et al.},
con $n_\text{tape}$ al posto di $\bar{n}_{c,i}$.

\subsection{Verifica interna: $h$ noto vs $h$ ricostruito}

Nel vostro caso $h$ è noto dalla costruzione del campione. La formula
\eqref{eq:spessore} fornisce quindi un \textbf{test di auto-consistenza}:

\[
  \varepsilon_h = h_\text{reco} - h_\text{noto}
  = \frac{\lambda}{2\pi}\cdot\frac{\varphi^{(k_1)}-\varphi^{(k_2)}}{\delta n}
  - n\cdot h_\text{layer}
\]

Se $|\varepsilon_h|$ è confrontabile con l'incertezza del calibro
($\sim 5\ \mu\text{m}$), il metodo è validato. Se è sistematicamente
maggiore, vi è un errore nella procedura di misura.

% ================================================================
\section{Analisi globale: fit lineare su tutte le coppie}
% ================================================================

\subsection{Impostazione del problema}

Per un gradino $n$ fissato, con $K$ misure di fase alle $K$ concentrazioni
$\{n_m^{(k)}\}_{k=1}^K$, il modello lineare \eqref{eq:retta} in forma
matriciale è:

\[
\underbrace{
\begin{pmatrix} \varphi_n^{(1)} \\ \varphi_n^{(2)} \\ \vdots \\ \varphi_n^{(K)} \end{pmatrix}
}_{\boldsymbol{\varphi}}
=
\underbrace{
\begin{pmatrix} 1 & -n_m^{(1)} \\ 1 & -n_m^{(2)} \\ \vdots & \vdots \\ 1 & -n_m^{(K)} \end{pmatrix}
}_{\mathbf{A}}
\underbrace{
\begin{pmatrix} B\cdot n_\text{tape} \\ B \end{pmatrix}
}_{\mathbf{x}},
\qquad B = \frac{2\pi\,n\,h_\text{layer}}{\lambda}.
\]

\subsection{Soluzione ai minimi quadrati}

Per $K > 2$ il sistema è \textbf{sovradeterminato}: si minimizza la somma dei
quadrati dei residui, $\min_{\mathbf{x}} \|\mathbf{A}\mathbf{x} - \boldsymbol{\varphi}\|^2$,
con soluzione:
\[
  \hat{\mathbf{x}} = \bigl(\mathbf{A}^\top \mathbf{A}\bigr)^{-1}
  \mathbf{A}^\top \boldsymbol{\varphi}.
\]

Le stime dei parametri fisici sono:
\begin{align*}
  \hat{n}_\text{tape} &= \frac{\hat{x}_1}{\hat{x}_2}, \\[4pt]
  \hat{h}_\text{layer} &= \frac{\lambda}{2\pi\,n}\,\hat{x}_2
  \quad (\text{se }\lambda\text{ è noto}).
\end{align*}



% ================================================================
\section{Propagazione degli errori}
% ================================================================

\subsection{Errore su $n_\text{tape}$ dalla formula \eqref{eq:indice}}

Indicando con $\sigma_\varphi$ l'incertezza sulla fase (dominata dalla
stabilità temporale del DHM, tipicamente $\sigma_\varphi \approx 0.2^\circ
\approx 3.5 \times 10^{-3}\ \text{rad}$ dopo mediazione spaziale) e con
$\sigma_{\delta n}$ l'incertezza su $\delta n$ (dalla precisione del
rifrattometro, tipicamente $\sigma_n \approx 2\times10^{-4}$):

\[
  \sigma_{n_\text{tape}} \approx
  \sqrt{
    \left(\frac{\partial n_\text{tape}}{\partial \varphi^{(k_1)}}\right)^2 \sigma_\varphi^2
    + \left(\frac{\partial n_\text{tape}}{\partial \varphi^{(k_2)}}\right)^2 \sigma_\varphi^2
    + \left(\frac{\partial n_\text{tape}}{\partial \delta n}\right)^2 \sigma_{\delta n}^2
  }
\]

dove le derivate parziali dalla formula \eqref{eq:indice} sono:
\begin{align*}
  \frac{\partial n_\text{tape}}{\partial \varphi^{(k_1)}}
  &= \frac{-\delta n\, \varphi^{(k_2)}}{\bigl(\varphi^{(k_1)}-\varphi^{(k_2)}\bigr)^2},\\
  \frac{\partial n_\text{tape}}{\partial \varphi^{(k_2)}}
  &= \frac{\delta n\, \varphi^{(k_1)}}{\bigl(\varphi^{(k_1)}-\varphi^{(k_2)}\bigr)^2},\\
  \frac{\partial n_\text{tape}}{\partial \delta n}
  &= \frac{\varphi^{(k_1)}}{\varphi^{(k_1)}-\varphi^{(k_2)}}.
\end{align*}

\paragraph{Osservazione importante.}
L'errore \textit{diverge} quando $\varphi^{(k_1)} \approx \varphi^{(k_2)}$,
cioè quando $\delta n \approx 0$. Le coppie di soluzioni con indici molto
simili (es.\ 60\% e 80\% se siete vicini al matching) danno stime più incerte.
Al contrario, coppie con $\delta n$ grande (es.\ 40\% e 80\%) sono più precise
per $n_\text{tape}$ ma richiedono attenzione all'unwrapping.

\subsection{Effetto della temperatura su $n_m^{(k)}$}

Il coefficiente termoottico delle soluzioni di saccarosio vale:
\[
  \frac{dn_m}{dT} \approx -1.5 \times 10^{-4}\ ^\circ\text{C}^{-1}.
\]
Una variazione di temperatura di $3\,^\circ\text{C}$ tra una misura e l'altra
introduce un errore sistematico $\Delta n_m \approx 4.5 \times 10^{-4}$,
confrontabile con la precisione del rifrattometro. È quindi necessario:
\begin{itemize}
  \item misurare $n_m^{(k)}$ con il rifrattometro immediatamente prima o dopo
        ogni acquisizione DHM (non usare valori tabulati),
  \item annotare la temperatura ad ogni misura.
\end{itemize}

% ================================================================
\section{Schema operativo dell'analisi}
% ================================================================

\begin{enumerate}
  \item \textbf{Per ogni acquisizione $(k)$}:
  \begin{enumerate}
    \item Verificare che la fase sia \textit{unwrapped} e priva di salti di
          $2\pi$ (ispezione visiva + test di monotonia tra gradini).
    \item Applicare la maschera su ciascun gradino $n = 1,\dots,8$.
    \item Calcolare la fase media $\langle\varphi_n^{(k)}\rangle$ e la
          deviazione standard $\sigma_n^{(k)}$ (misura della variabilità
          spaziale dello spessore nel gradino).
  \end{enumerate}

  \item \textbf{Analisi per coppie} $(k_1, k_2)$:
  \begin{enumerate}
    \item Calcolare $\delta n = n_m^{(k_1)} - n_m^{(k_2)}$ dai valori del
          rifrattometro.
    \item Per ogni gradino $n$, calcolare:
    \[
      n_\text{tape}^{(n)} = \frac{\delta n \cdot \langle\varphi_n^{(k_1)}\rangle}
      {\langle\varphi_n^{(k_1)}\rangle - \langle\varphi_n^{(k_2)}\rangle}
      + n_m^{(k_1)},
    \]
    \[
      h_\text{reco}^{(n)} = \frac{\lambda}{2\pi}\cdot
      \frac{\langle\varphi_n^{(k_1)}\rangle - \langle\varphi_n^{(k_2)}\rangle}{\delta n}.
    \]
    \item Verificare $h_\text{reco}^{(n)} \approx n \cdot h_\text{layer}$.
    \item Verificare che $n_\text{tape}^{(n)}$ sia consistente al variare di $n$.
  \end{enumerate}

  \item \textbf{Analisi globale} (per gradino $n$ fissato):
  \begin{enumerate}
    \item Tracciare $\langle\varphi_n^{(k)}\rangle$ vs $n_m^{(k)}$
          per tutti i valori $k$.
    \item Eseguire un fit lineare (minimi quadrati pesati con $1/\sigma_n^{(k)}$).
    \item Estrarre $n_\text{tape}$ dallo zero della retta interpolata.
    \item Verificare che le rette di gradini diversi abbiano zeri coincidenti.
  \end{enumerate}
\end{enumerate}

% ================================================================
\section{Riepilogo: differenze rispetto al paper originale}
% ================================================================

\begin{center}
\begin{tabular}{p{4.5cm}p{4.5cm}p{4.5cm}}
\hline\hline
& \textbf{Paper (cellula)} & \textbf{Questo esperimento} \\
\hline
Campione & Vivo, eterogeneo, mobile &
  Solido, omogeneo, fisso \\[6pt]
Incognite & $h_i$ e $\bar{n}_{c,i}$ per ogni pixel &
  $n_\text{tape}$ sola (globale) \\[6pt]
Spessore & Ignoto & Noto: $h_n = n\cdot h_\text{layer}$ \\[6pt]
N.\ soluzioni & 2 (sistema esattamente determinato) &
  4 (sistema sovradeterminato) \\[6pt]
Analisi & Pixel per pixel & Per gradino (media spaziale) \\[6pt]
Verifica & Non possibile & $h_\text{reco} \stackrel{?}{=} h_\text{noto}$ \\[6pt]
Phase wrapping & Meno critico ($h \sim \mu\text{m}$) &
  \textbf{Critico} ($h \sim 30\ \mu\text{m}$, più giri) \\[6pt]
$\lambda$ richiesto & Sì, per le formule & Per $h$; non per $n_\text{tape}$ \\
\hline\hline
\end{tabular}
\end{center}

% ================================================================
\section{Questioni aperte}
% ================================================================

Le seguenti informazioni non erano disponibili al momento della stesura
e influenzano la precisione quantitativa dell'analisi:

\begin{enumerate}
  \item \textbf{Lunghezza d'onda $\lambda$}: necessaria per convertire la
        pendenza del fit in $h_\text{layer}$ e per stimare il numero di giri
        di fase. Da ricavare dalla scheda tecnica del DHM.
  \item \textbf{Quarta concentrazione}: chiarire se ``100\%'' è acqua pura,
        soluzione satura, o concentrazione in g/L.
  \item \textbf{Qualità dell'unwrapping}: verificare esplicitamente la
        monotonia della fase tra gradini per ogni acquisizione.
  \item \textbf{Spessore indipendente}: se possibile, misurare $h_\text{layer}$
        con un metodo più preciso del calibro (profilometro, microscopia confocale).
\end{enumerate}



% ! ================================================================
% ! METTERE IN APPENDICE QUANTO SEGUE:
\pagecolor{white}
{\color{green!30!black}
\section{Fase interferometrica e phase unwrapping in Digital Holographic Microscopy}

Nel Digital Holographic Microscopy (DHM), l’informazione principale ricavata dalla ricostruzione olografica è la fase ottica del campo complesso, che codifica lo spessore ottico del campione. La fase misurata è direttamente proporzionale alla differenza di cammino ottico tra il raggio che attraversa il campione e quello propagante nel mezzo circostante:

\begin{equation}
\varphi(x,y) = \frac{2\pi}{\lambda}\,\Delta n(x,y)\,h(x,y),
\end{equation}

dove $\lambda$ è la lunghezza d’onda di illuminazione, $\Delta n = n_{\text{campione}} - n_{\text{mezzo}}$ è la differenza di indice di rifrazione e $h(x,y)$ rappresenta lo spessore locale del campione. La fase $\varphi$ è una grandezza angolare non limitata superiormente e può assumere valori arbitrariamente grandi in presenza di campioni spessi o fortemente contrastati in indice di rifrazione.

\subsection{Fase wrapped e periodicità della misura}

Il sensore (tipicamente CCD o CMOS) non misura direttamente la fase $\varphi$, ma registra l’intensità del pattern interferometrico, che dipende dalle quantità $\cos(\varphi)$ e $\sin(\varphi)$. Poiché tali funzioni sono periodiche con periodo $2\pi$, la fase recuperata numericamente è definita solo modulo $2\pi$:

\begin{equation}
\varphi_{\text{wrapped}} \in [-\pi, \pi].
\end{equation}

Di conseguenza, la misura non consente di distinguere tra $\varphi$, $\varphi + 2\pi k$ con $k \in \mathbb{Z}$. Questa ambiguità è nota come \textit{phase wrapping}.

Una rappresentazione utile consiste nel considerare la fase wrapped come una funzione a valori periodici, che introduce discontinuità artificiali ogni volta che la fase reale supera i limiti dell’intervallo $[-\pi,\pi]$. In tali condizioni, la fase osservata presenta un andamento a “denti di sega”, pur derivando da una distribuzione fisicamente continua.

\subsection{Interpretazione spaziale e frange di interferenza}

Nel caso di un campione con variazione graduale di spessore, la fase fisica varia in modo continuo nello spazio. Tuttavia, la rappresentazione wrapped introduce discontinuità apparenti che si manifestano come frange interferometriche.

È importante sottolineare che tali frange non corrispondono necessariamente a discontinuità fisiche del campione, ma rappresentano isofasi della quantità $\varphi_{\text{wrapped}}$. La loro densità spaziale è legata al gradiente di fase e quindi allo spessore ottico locale.

Nel caso di discontinuità geometriche marcate, come gradini micrometrici, il salto di fase può risultare molto elevato. Ad esempio, per un salto di spessore $h \sim 30\,\mu m$ in una soluzione con $\Delta n \sim 0.08$, si ottiene:

\begin{equation}
\Delta \varphi \approx \frac{2\pi}{\lambda} \Delta n \, h,
\end{equation}

che, per $\lambda \sim 650\,nm$, può raggiungere valori dell’ordine di decine di radianti. In tali condizioni, la variazione di fase tra campioni adiacenti può eccedere la soglia di $\pi$, compromettendo la continuità locale della fase wrapped.

\subsection{Phase unwrapping}

L’operazione di \textit{phase unwrapping} consiste nel ricostruire una stima continua della fase $\varphi$ a partire dalla sua versione wrapped, mediante l’aggiunta di multipli interi di $2\pi$:

\begin{equation}
\varphi(x,y) = \varphi_{\text{wrapped}}(x,y) + 2\pi k(x,y),
\end{equation}

dove $k(x,y) \in \mathbb{Z}$ è determinato numericamente.

Gli algoritmi di unwrapping si basano generalmente sull’assunzione che la differenza di fase tra pixel adiacenti soddisfi la condizione:

\begin{equation}
|\varphi(x_i,y_i) - \varphi(x_j,y_j)| < \pi.
\end{equation}

Questa ipotesi garantisce la ricostruzione univoca della fase in presenza di variazioni spaziali sufficientemente lente. Tuttavia, in presenza di gradienti elevati o discontinuità geometriche, tale condizione può essere violata, rendendo il problema ambiguo.

In tali casi, l’unwrapping può risultare sensibile alla scelta dell’algoritmo e alla presenza di rumore, con possibile propagazione locale degli errori, soprattutto nei metodi basati su integrazione spaziale diretta.

\subsection{Osservazioni finali}

La fase unwrapped rappresenta la quantità fisicamente significativa per la ricostruzione quantitativa dello spessore ottico e deve essere utilizzata per qualsiasi stima metrologica del campione. Al contrario, la fase wrapped è utile principalmente per l’analisi qualitativa delle frange interferometriche.

In applicazioni con forti discontinuità spaziali, è fondamentale verificare la correttezza dell’unwrapping mediante controlli di consistenza, in quanto errori locali possono propagarsi e influenzare la stima globale dello spessore.
}